\documentclass[french]{llncs}

\usepackage[english, french]{babel}
\usepackage[utf8]{inputenc}

\usepackage{graphicx}
\usepackage{amsmath}

\usepackage{listings}
\usepackage{subfig}
\usepackage{hyperref}
\usepackage{booktabs}

%-----------------------------------------
\title{Structure topologique ou caractéristiques individuelles~? Prédire le pays d'un aéroport dans un réseau aérien}

\author{ Yennipena Jérémie \textsc{Coulibaly} \\
         Michel Junior \textsc{Delanoue} }

\institute{Université Claude Bernard Lyon 1, France \\
\url{https://github.com/cyjere/rendu-tp-gnn-Coulibaly-and-Delanoue}}

\lstset{
    basicstyle=\small,
    captionpos=b
}

\providecommand{\motcles}[1]{\textbf{Mots-clés:} #1}

\pagestyle{headings}
%-----------------------------------------

\begin{document}

    \renewcommand{\thelstlisting}{\arabic{lstlisting}}

    \maketitle

    \selectlanguage{english}
    \begin{abstract}
        We study node classification on a worldwide airport network (3\,363 airports, 13\,547 routes) where the task is to 
        predict an airport's country from its attributes (longitude, latitude, population) and/or the graph of flight connections.
         We compare classical machine learning and deep learning baselines (Logistic Regression, Random Forest, MLP) 
         trained on node features alone against graph neural networks (GCN, GAT, SGC) trained either on the graph structure 
         alone or on features combined with structure, all evaluated with class-weighted losses and macro-F1 over five random 
         splits to handle the severe class imbalance (212 countries, from 1 to 650 airports each). Contrary to the usual 
         expectation, adding graph structure never improves over the best feature-only model 
         (Random Forest, macro-F1 $0.779 \pm 0.018$); structure alone is no better than chance, 
         and combining it with features consistently hurts GCN/GAT/SGC relative to Random Forest. 
         We show this is explained by low edge homophily (63.2\% of routes connect two airports of the same country) 
         concentrated on the busiest international hubs, which dilutes the geographic signal that message passing relies on.

         \begin{keywords}
        graph neural networks, node classification, class imbalance, airport network, ablation study
    \end{keywords}
    \end{abstract}

    \selectlanguage{french}

%-----------------------------------------

    \section{Introduction} \label{intro}

    Les réseaux aériens forment un graphe naturel~: les aéroports sont des nœuds, les liaisons directes des arêtes. 
    On peut se demander si la \emph{structure} de ce graphe apporte une information utile pour prédire un attribut d'un aéroport,
     au-delà de ses seules caractéristiques individuelles. Nous étudions cette question sur la tâche de prédiction du 
     \textbf{pays} d'un aéroport, à partir d'un graphe de 3\,363 aéroports et 13\,547 liaisons, chaque nœud portant sa longitude,
      sa latitude, la population de la ville desservie et son pays (variable cible).

    Notre problématique est la suivante~: 
    \emph{dans quelle mesure la structure topologique du réseau aérien améliore-t-elle la prédiction du pays d'un aéroport par 
    rapport à ses seules caractéristiques individuelles~?} Pour y répondre, nous comparons trois conditions expérimentales~: 
    (A) des modèles classiques de machine learning et de deep learning entraînés uniquement sur les caractéristiques des nœuds~; 
    (B) des GNN entraînés uniquement sur la structure du graphe, sans aucune caractéristique~; 
    (C) les mêmes GNN entraînés sur les caractéristiques \emph{et} la structure. Cette décomposition en trois conditions, 
    permet d'isoler la contribution propre de la topologie.

    \section{État de l'art} \label{sect:eda}

    Pour la condition graphe, nous retenons trois architectures. Le \textbf{GCN} \cite{Kipf2017} est la référence historique~: 
    chaque couche moyenne la représentation d'un nœud avec celle de ses voisins, pondérée par le degré. Le \textbf{SGC}
    (\emph{Simplifying Graph Convolution}) \cite{Wu2019} simplifie le GCN en supprimant les transformations non linéaires
    intermédiaires et les matrices de poids entre les couches, afin de ne conserver essentiellement que la propagation de 
    l'information dans le graphe. Comparé au GCN, il permet de vérifier si cette non-linéarité apporte réellement quelque chose. 
    Le \textbf{GAT} \cite{Velickovic2018} remplace la moyenne pondérée par le degré par un poids d'attention appris par paire 
    de nœuds~: nous le retenons pour tester si un mécanisme capable d'ignorer des voisins non pertinents peut compenser 
    un éventuel manque d'homophilie du graphe.

    Pour les baselines sans graphe, le \textbf{Random Forest} \cite{Breiman2001} est un classifieur non linéaire robuste, 
    bien adapté à de petites features numériques continues comme des coordonnées géographiques. 
    Le \textbf{perceptron multicouche} \cite{Rumelhart1986} sert de baseline deep learning sans information structurelle. 
    Enfin, comme le jeu de données présente un déséquilibre de classes sévère (212 pays, de 1 à 650 aéroports), 
    nous nous appuyons sur la revue de \cite{Wang2021} sur la classification de données déséquilibrées pour justifier 
    le choix d'une pondération de la perte par classe et de métriques moyennées par classe (macro-F1) 
    plutôt que l'accuracy brute.

    \section{Données et protocole} \label{sect:protocol}

    \paragraph{Données.} Le graphe est chargé depuis un fichier GraphML~: 3\,363 nœuds (aéroports), 13\,547 arêtes (liaisons), 
    avec pour chaque nœud \texttt{lon}, \texttt{lat}, \texttt{population}, \texttt{city\_name} et \texttt{country} (cible). 
    La cible comporte 212 classes brutes très déséquilibrées (USA~: 650 aéroports, médiane~: 4, 62 pays à un seul aéroport), 
    ce qui rend impossible toute évaluation sur les plus petites classes.

    \paragraph{Filtrage et split.} Nous retirons les pays à moins de 3 aéroports (\texttt{MIN\_CLASS\_SIZE=3}), ramenant 
    le problème à 126 classes sur 3\,253 nœuds et 12\,272 arêtes. Un split stratifié classique échoue encore sur ce seuil~: 
    après une première coupe train/reste, de nombreuses classes à 3 éléments n'ont plus qu'un seul exemple, 
    impossible à répartir entre validation et test. Nous utilisons donc un split par classe garantissant au moins 
    un exemple dans chacun des trois ensembles, avec un repli 
    proportionnel pour les classes plus grandes~: 69,0\,\%~/~15,5\,\%~/~15,5\,\%, répété sur 5 seeds.

    \paragraph{Features.} Les coordonnées brutes \texttt{lon/lat} posent deux problèmes géométriques~: 
    la longitude est circulaire (deux points proches peuvent être à $\pm180°$) et les pôles sont des singularités. 
    Nous testons donc un encodage sur la sphère unité 
    $(x,y,z) = (\cos(\text{lat})\cos(\text{lon}), \cos(\text{lat})\sin(\text{lon}), \sin(\text{lat}))$, combiné à $\log(1+\text{population})$, 
    standardisé sur le train uniquement. Cet encodage améliore tous les modèles testés (section~\ref{sect:results}) et devient 
    la représentation par défaut.

    \paragraph{Déséquilibre et hyperparamètres.} Les modèles sklearn utilisent \texttt{class\_weight="balanced"}~; 
    les réseaux de neurones (MLP, GCN, GAT, SGC) utilisent une \texttt{CrossEntropyLoss} pondérée par l'inverse de la fréquence 
    de classe, calculée sur le train de chaque seed. 
    L'architecture du MLP (256 unités, 2 couches, \texttt{lr=0,005}, \texttt{dropout=0,1}) a été choisie par une recherche sur 
    la validation du seed~0 uniquement (macro-F1~: 0,593~$\rightarrow$~0,684), jamais sur le test.

    \newpage
    \bibliographystyle{splncs03}
    \bibliography{bibM2.bib}

\end{document}
